\section{Hyperparameter Details}
\label{app:hyperparameter-details}
Table~\ref{tab:rpc-hyperparameters} reports the \model hyperparameters used in the main experiments. The parameters are the rarity exponent $\alpha$, frequency-penalty weight $\lambda$, and soft class-balancing weight $w$, matching the notation in Eq.~\ref{eq:class-balance-score}.

\begin{table}[h]
\centering
\caption{\model hyperparameters used in the main experiments.}
\label{tab:rpc-hyperparameters}
\small
\setlength{\tabcolsep}{5pt}
\begin{tabular}{l c c c}
\toprule
\rowcolor{sectiongray} \textbf{Dataset} & $\boldsymbol{\alpha}$ & $\boldsymbol{\lambda}$ & \textbf{w} \\
\midrule
ImageNet10 & 0.03 & 0.5 & 15 \\
CelebA & 0.03 & 0.5 & 15 \\
COCO & 0.03 & 0.5 & 15 \\
KITTI & 0.01 & 5.0 & 15 \\
\bottomrule
\end{tabular}
\end{table}

The larger frequency-penalty weight used for KITTI is consistent with the relative size of the generated repair set. All four settings add 1k generated images, but KITTI has only 5,000 source training images, compared with 10,400 for ImageNet10, 20,000 for COCO, and 80,000 for CelebA. The generated images therefore form a larger fraction of the KITTI repair candidate pool, making a stronger penalty on already-supplied attribute values useful for preserving source-data complementarity. This dataset-dependent shift supports the design of the score: the frequency-penalty term adapts the selector to the amount of generated repair data relative to the available source pool.

Table~\ref{tab:kitti-rpc-sensitivity} gives the KITTI sensitivity across source-selection ratios. Across several \model parameter settings, the test mAP50--95 changes only modestly, suggesting that the selector is not brittle to small changes in $\alpha$, $\lambda$, and $w$ once the generated-data frequency penalty is present.

\begin{table}[h]
\centering
\caption{KITTI \model sensitivity. Metric: test mAP50--95.}
\label{tab:kitti-rpc-sensitivity}
\small
\setlength{\tabcolsep}{4pt}
\begin{tabular}{c c c c c c c c}
\toprule
\rowcolor{sectiongray} $\boldsymbol{\alpha}$ & $\boldsymbol{\lambda}$ & \textbf{w} & \textbf{0.1} & \textbf{0.3} & \textbf{0.5} & \textbf{0.7} & \textbf{0.9} \\
\midrule
0.01 & 5.0 & 15 & \textbf{0.5372} & \textbf{0.5646} & \textbf{0.5773} & 0.5780 & \textbf{0.5850} \\
0.03 & 0.5 & 15 & 0.5358 & 0.5637 & 0.5764 & 0.5774 & 0.5842 \\
0.03 & 5.0 & 8 & 0.5366 & 0.5641 & 0.5769 & \textbf{0.5784} & 0.5847 \\
0.10 & 5.0 & 15 & 0.5349 & 0.5640 & 0.5768 & 0.5776 & 0.5839 \\
0.80 & 10.0 & 4 & 0.5337 & 0.5629 & 0.5756 & 0.5769 & 0.5828 \\
\bottomrule
\end{tabular}
\end{table}

Table~\ref{tab:coco-rpc-sensitivity} shows the same pattern on COCO pose-estimation repair. The selected main setting remains strongest for most ratios, while nearby \model parameter settings stay within a narrow AP band.

\begin{table}[h]
\centering
\caption{COCO \model sensitivity. Metric: COCO keypoint AP.}
\label{tab:coco-rpc-sensitivity}
\small
\setlength{\tabcolsep}{4pt}
\begin{tabular}{c c c c c c c c}
\toprule
\rowcolor{sectiongray} $\boldsymbol{\alpha}$ & $\boldsymbol{\lambda}$ & \textbf{w} & \textbf{0.1} & \textbf{0.3} & \textbf{0.5} & \textbf{0.7} & \textbf{0.9} \\
\midrule
0.03 & 0.5 & 15 & \textbf{0.6183} & \textbf{0.6306} & \textbf{0.6317} & 0.6319 & \textbf{0.6355} \\
0.01 & 0.5 & 15 & 0.6152 & 0.6298 & 0.6309 & 0.6308 & 0.6337 \\
0.03 & 0.5 & 8 & 0.6165 & 0.6289 & 0.6305 & \textbf{0.6323} & 0.6344 \\
0.03 & 0.5 & 0 & 0.6148 & 0.6276 & 0.6292 & 0.6306 & 0.6328 \\
0.00 & 0.0 & 15 & 0.6042 & 0.6288 & 0.6203 & 0.6313 & 0.6268 \\
\bottomrule
\end{tabular}
\end{table}

